%% commands.tex
%% Custom macros and theorem environments for the BNP thesis.

% Draft annotation: \MC{note} renders as coloured inline text
\newcommand{\MC}[1]{\textcolor{red}{\textbf{[MC: #1]}}}

% ============================================================
%  THEOREM ENVIRONMENTS
%  Numbered per section (chapter.section.number).
% ============================================================
\theoremstyle{plain}
\newtheorem{theorem}{Theorem}[section]
\newtheorem{proposition}[theorem]{Proposition}
\newtheorem{lemma}[theorem]{Lemma}
\newtheorem{corollary}[theorem]{Corollary}

\theoremstyle{definition}
\newtheorem{definition}[theorem]{Definition}
\newtheorem{assumption}[theorem]{Assumption}
\newtheorem{example}[theorem]{Example}
\newtheorem{remark}[theorem]{Remark}
\newtheorem{ttheorem}[theorem]{Tentative Theorem}
\newtheorem*{prop}{Proposition}
\newtheorem*{cor}{Corollary}

% ============================================================
%  Bayesian Nonparametric Toolbox
% ============================================================

% For spaces
\newcommand{\X}{\mathbb{X}}
\newcommand{\Y}{\mathbb{Y}}
\newcommand{\probx}{\mathcal{P}(\mathbb{X})}
\newcommand{\probspace}{\left( \Omega, \mathscr{F}, \mathbb{P} \right)}
\newcommand{\probprobx}{\mathcal{P}(\mathcal{P}(\X))}
\newcommand{\mbx}{\mathcal{M}_B(\X)}
\newcommand{\mx}{\mathcal{M}(\X)}
\newcommand{\probmb}{\mathcal{P}(\mbx)}
\newcommand{\lipR}{\mathrm{Lip}_1(\R,\R)}
\newcommand{\lipX}{\mathrm{Lip}_1(\X,\R)}
\newcommand{\lipK}{\mathrm{Lip}_1([-K,K],\R)}
\newcommand{\lipXstar}{\mathrm{Lip}_1^*(\X,\R)}
\newcommand{\lipKzero}{\mathrm{Lip}_1^0([-K,K],\R)}
\newcommand{\dstar}{\D^*}
\newcommand{\ellstar}{\ell^\infty(\dstar)}

% For objects
\newcommand{\mustar}{\widetilde{\mu}^*}

% For distances
\newcommand{\bl}{\mathrm{BL}}
\newcommand{\ad}{\mathrm{AD}}
\newcommand{\wext}{\W_*}
\newcommand{\wbl}{\W_{\mathrm{BL}}}
\newcommand{\wad}{\W_{\mathrm{AD}}}
\newcommand{\wtr}{\W_{d_{\X,t}}}
\newcommand{\dlip}{d_{\mathrm{Lip}}}
\newcommand{\WW}{\mathcal{W}_{\mathcal{W}}}








% ============================================================
% Some basic notations
% ============================================================
\newcommand{\R}{\mathbb{R}}
\newcommand{\E}{\mathbb{E}}
\newcommand{\N}{\mathcal{N}}          % Poisson random measure / calligraphic N
\newcommand{\M}{\mathcal{M}}
\newcommand{\G}{\mathcal{G}}
\newcommand{\F}{\mathcal{F}}
\newcommand{\D}{\mathcal{D}}
\newcommand{\W}{\mathcal{W}}
\newcommand{\T}{\mathcal{T}}
\newcommand{\law}{\mathcal{L}}
\newcommand{\var}{\mathrm{Var}}
\newcommand{\cov}{\mathrm{Cov}}
\newcommand{\iid}{\mathrm{i.i.d}}

